\section{PLANTEAMIENTO DEL PROBLEMA}

\subsection{Antecedentes}

La persistencia y transmisión de datos son procesos fundamentales en el desarrollo de software, donde la elección del formato de serialización impacta directamente la eficiencia general del sistema. Los archivos constituyen los elementos base para la persistencia y transferencia de información, por lo que comprender su comportamiento en memoria es clave para optimizar estos procesos técnicos.

El rendimiento en la manipulación de archivos abarca dimensiones como el tiempo de lectura, tiempo de instanciación, consumo de memoria principal y uso de CPU. En la práctica, se suele ignorar cómo la elección de un formato afecta la eficiencia interna del programa, especialmente cuando este se ejecuta bajo recursos de hardware restringidos. Realizar un análisis comparativo y cuantitativo para la selección objetiva del tipo de archivo resulta esencial para garantizar la estabilidad de las aplicaciones.

Actualmente, el uso de los formatos CSV, JSON y XML se encuentra ampliamente extendido en diferentes contextos de almacenamiento e intercambio de datos, siendo objeto recurrente de análisis de rendimiento \cite{ieee_icoailo_2025, ieee_les_2024}. A pesar de la literatura existente, se requieren evidencias empíricas en escenarios concretos que evalúen el sobrecosto real de estas tecnologías bajo modelos de carga específicos.

La falta de claridad sobre la huella computacional de cada formato dificulta la anticipación de cuellos de botella al procesar grandes volúmenes de información en memoria principal. En archivos con estructuras sintácticamente complejas o niveles de anidamiento profundos, el proceso de parseo puede elevar el uso de CPU y degradar el rendimiento global del sistema \cite{ieee_parsing_2008}.

De igual forma, la ineficiencia en el manejo de memoria RAM se manifiesta como un problema crítico derivado del modelo de representación utilizado. La carga completa en memoria mediante árboles DOM en documentos extensos multiplica la huella de memoria (\textit{memory footprint}), pudiendo sobrepasar la capacidad asignada a la Máquina Virtual de Java (JVM) o degradar el tiempo de respuesta \cite{dke_memory_2006, hoxmeier2000system}.

La realización temprana de pruebas de rendimiento y la caracterización de sobrecostos sintácticos permiten orientar la selección de formatos y mitigar riesgos en las decisiones de arquitectura de software \cite{smith2002performance}.


\subsection{Descripción de la situación problemática}

Como se documentó en los antecedentes, los formatos CSV, JSON y XML se emplean masivamente para almacenar e intercambiar información. Sin embargo, su procesamiento mediante modelos de carga completa en memoria (DOM) suele generar sobrecostos no cuantificados y un elevado uso de recursos en la JVM al manipular archivos de gran tamaño.

Actualmente no se dispone de una caracterización experimental y comparativa sobre el rendimiento real de estos tres formatos —específicamente en tiempos de procesamiento y consumo de memoria RAM (pico y retenida)— bajo entornos computacionales controlados y estandarizados. Esta carencia de evidencia técnica obliga a desarrolladores y arquitectos de software a seleccionar los formatos de archivo e interfaces de parseo con base en la costumbre o la intuición, desconociendo el impacto físico exacto sobre el hardware.


\subsection{Formulación del problema}

A partir del planteamiento anterior, la pregunta de investigación que orienta este estudio es:

¿Qué efecto tiene el procesamiento mediante la carga completa en memoria utilizando el modelo basado en árbol (DOM) (variable independiente) sobre los tiempos de ejecución y el consumo de memoria RAM (variables dependientes) al parsear e inspeccionar archivos de formato CSV, JSON y XML con volúmenes de datos controlados en una arquitectura computacional estandarizada?

Esta pregunta se sistematiza en las siguientes sub-preguntas:
\begin{enumerate}
    \item[a)] ¿Cómo influye la estructura sintáctica del modelo DOM en el tiempo de procesamiento (lectura del archivo y construcción del árbol en memoria) de archivos CSV, JSON y XML bajo volúmenes de datos equivalentes?
    \item[b)] ¿Qué efecto produce la sobrecarga de instanciación de nodos del modelo DOM en el consumo de memoria RAM según la sintaxis y el nivel de anidamiento de cada formato?
    \item[c)] ¿Existen diferencias en el tamaño y la estructura inicial de los archivos CSV, JSON y XML que deban medirse primero para no confundir el peso de sus etiquetas con fallos de rendimiento?
\end{enumerate}

El proyecto se aborda mediante un diseño experimental cuantitativo bajo un entorno computacional de pruebas estandarizado, lo que permite aislar las variables intervinientes y determinar la causalidad directa entre la sintaxis del formato y el impacto físico en el hardware.


\subsection{Consecuencias de no investigar}

Si la ausencia de evidencia técnica persiste, las decisiones de selección de formatos continuarán tomándose de forma empírica, incrementando el riesgo de desarrollar aplicaciones vulnerables a la saturación de memoria RAM y tiempos de procesamiento ineficientes. Cada sistema diseñado bajo esta incertidumbre queda expuesto a fallos de ejecución o refactorizaciones costosas en etapas avanzadas de producción que podrían haberse prevenido desde la fase de arquitectura.


\subsection{Delimitación}

Este estudio se delimita a medir y comparar el rendimiento computacional —tiempos de lectura e instanciación en memoria y consumo de memoria RAM— de los formatos de archivo CSV, JSON y XML, evaluando su procesamiento mediante la carga completa en memoria utilizando el modelo basado en árbol (DOM), sobre la Máquina Virtual de Java (JVM) en una plataforma de hardware con características controladas.

Queda fuera del alcance la velocidad de internet o de red como factor explicativo, así como la evaluación de procesamiento concurrente o multi-hilo, el uso de técnicas de procesamiento por flujos (\textit{streaming}), la interacción con sistemas de bases de datos y el análisis de aspectos de seguridad o encriptación de la información.